% Input-safe canonical proof module.
% Host requirements: amsmath, amssymb, amsthm, mathtools, mathrsfs,
% array, booktabs, enumitem, hyperref, and theorem environments
% theorem, proposition, lemma, and counterexample.
\begingroup
\def\J{\mathcal J_2}
\def\V{\mathscr V}
\def\Psd{\mathcal P_2}
\def\Ppos{\mathcal P_2^{++}}
\def\Nfuture{\mathcal N_+^\times}
\def\rankC{\operatorname{rank}_{\mathbb C}}
\def\dd{\mathrm d}
\def\SpanR{\operatorname{span}_{\mathbb R}}

\section[Exact presymplectic reduction of the SU(3)-Gram worldline model]
{Exact presymplectic reduction of the\\
\texorpdfstring{$SU(3)$}{SU(3)}-Gram worldline model}
\label{wpr:sec:worldline-reduction}

The constrained Hamiltonian framework used below goes back to Dirac's
analysis of singular Lagrangians and constrained Hamiltonian systems
\cite{Dirac1950}.  The separation of constraint surface, characteristic
distribution, and reduced phase space is the finite-dimensional
presymplectic constraint problem treated systematically by
Gotay--Nester--Hinds \cite{GotayNesterHinds1978}; the regular quotient
step is the symplectic reduction mechanism of Marsden--Weinstein
\cite{MarsdenWeinstein1974}.  These references supply the general
framework only.  Every constraint rank, orbit, stabilizer, kernel,
one-form, two-form, quotient map, sign, equation of motion, and singular
counterexample used for this particular (SU(3))-Gram model is computed
explicitly below.

\subsection{Objects, constants, and equations}
Let
\[
 \J=\{A\in\operatorname{Mat}_{2\times2}(\mathbb C):A=A^*\},
\qquad
 \langle A,B\rangle_L=
 \frac12\bigl(\operatorname{tr}A\operatorname{tr}B-\operatorname{tr}(AB)\bigr).
\]
Then \(\langle A,A\rangle_L=\det A\), and the pairing is nondegenerate.
Write
\[
 \Psd=\{p\in\J:p\succeq0\},\quad
 \Ppos=\{p\in\J:p\succ0\},\quad
 \Nfuture=\{p\in\Psd:\rankC p=1\}.
\]
Retain
\[
 \V=\operatorname{Mat}_{3\times2}(\mathbb C),\quad
 \alpha>0,\quad c>0,\quad m_0\geq0,
\]
and define
\begin{align}
 P_\alpha(Z)&=\alpha Z^*Z,\label{wpr:eq:P}\\
 D P_{\alpha,Z}(\delta Z)&=
 \alpha(\delta Z^*Z+Z^*\delta Z),\label{wpr:eq:DP}\\
 \Phi_{m_0,\alpha,c}(Z)&=
 \frac12\left(\det P_\alpha(Z)-m_0^2c^2\right).\label{wpr:eq:Phi}
\end{align}
The history configuration space on a retained parameter interval \(I\) is
\[
 \mathfrak C_I=
 C^\infty(I,\J)\times C^\infty(I,\V)\times C^\infty(I,\mathbb R),
\]
with coordinates \((x,Z,e)\).  Because the Lagrangian is first order
in \(x\), contains no \(\dot Z\), and contains no \(\dot e\), the
corresponding instantaneous presymplectic phase carrier is
\[
 \mathcal E={\J}_x\times{\V}_Z\times\mathbb R_e.
\]
Define
\begin{equation}
 \theta=\langle P_\alpha(Z),\dd x\rangle_L,\qquad
 \Omega=-\dd\theta,\qquad
 \mathcal H=e\Phi_{m_0,\alpha,c}.
 \label{wpr:eq:forms}
\end{equation}
For tangent vectors \(v=(v_x,v_Z,v_e)\) and \(w=(w_x,w_Z,w_e)\),
\begin{equation}
 \Omega(v,w)=
 \langle v_x,D P_{\alpha,Z}(w_Z)\rangle_L
 -\langle w_x,D P_{\alpha,Z}(v_Z)\rangle_L.
 \label{wpr:eq:Omega}
\end{equation}

\begin{proposition}[Exact presymplectic Euler--Lagrange equation]
\label{wpr:prop:EL}
A path \((x,Z,e)\) is stationary under compactly supported variations
if and only if
\begin{align}
 \det P_\alpha(Z)&=m_0^2c^2,\label{wpr:eq:EL1}\\
 \frac{\dd}{\dd\tau}P_\alpha(Z)&=0,\label{wpr:eq:EL2}\\
 \left\langle D P_{\alpha,Z}(\delta Z),
 \dot x-eP_\alpha(Z)\right\rangle_L&=0
 \quad\text{for every }\delta Z\in\V.\label{wpr:eq:EL3}
\end{align}
Equivalently, \(\iota_{\dot\gamma}\Omega=\dd\mathcal H\).
\end{proposition}
\begin{proof}
Since
\[
 D(\det)_p(A)=2\langle p,A\rangle_L,
\]
one has
\[
 \dd\mathcal H(w)=
 \Phi_{m_0,\alpha,c}w_e+
 e\langle P_\alpha(Z),D P_{\alpha,Z}(w_Z)\rangle_L.
\]
Insert \eqref{wpr:eq:Omega} into
\(\iota_{\dot\gamma}\Omega=\dd\mathcal H\) and test successively
\(w_x,w_e,w_Z\).  This gives \eqref{wpr:eq:EL2},
\eqref{wpr:eq:EL1}, and \eqref{wpr:eq:EL3}.
\end{proof}

\subsection{Colour quotient and orbit types}
The left action is \(U\cdot(x,Z,e)=(x,UZ,e)\).

\begin{lemma}[Local colour is a zero-moment null symmetry]
\label{wpr:lem:null}
For \(\xi\in\mathfrak{su}(3)\), let \(\xi^\#=(0,\xi Z,0)\).  Then
\[
 D P_{\alpha,Z}(\xi Z)=0,\qquad
 \iota_{\xi^\#}\theta=0,\qquad
 \iota_{\xi^\#}\Omega=0.
\]
Thus the zero map to \(\mathfrak{su}(3)^*\) is a presymplectic moment
map.  The action is invariant under every \(\tau\)-dependent
\(U:I\to SU(3)\), with no connection term.
\end{lemma}
\begin{proof}
Because \(\xi^*=-\xi\),
\[
 D P_{\alpha,Z}(\xi Z)=
 \alpha(Z^*\xi^*Z+Z^*\xi Z)=0.
\]
The form identities follow from \eqref{wpr:eq:forms} and
\eqref{wpr:eq:Omega}.  Pointwise,
\((UZ)^*(UZ)=Z^*Z\), even when \(U\) depends on \(\tau\).
\end{proof}

\begin{theorem}[Gram quotient, fibres, and stabilizers]
\label{wpr:thm:quotient}
The map \(P_\alpha\) induces a homeomorphism
\[
 \V/SU(3)\xrightarrow{\ \sim\ }\Psd,\qquad [Z]\longmapsto P_\alpha(Z).
\]
Its orbit types are
\begin{center}
\begin{tabular}{c c c}
\toprule
\(\rankC Z\)&\text{fibre}&\(\operatorname{Stab}_{SU(3)}(Z)\)\\
\midrule
\(2\)&\(SU(3)\)&\(\{I_3\}\)\\
\(1\)&\(SU(3)/SU(2)\cong S^5\)&\(SU(2)\)\\
\(0\)&a point&\(SU(3)\)\\
\bottomrule
\end{tabular}
\end{center}
The rank-two quotient is smooth and free.  Rank one is a smooth
nonprincipal orbit-type stratum with nontrivial stabilizer; the zero
vertex is the singular apex of the total stratified quotient.
\end{theorem}
\begin{proof}
If \(Z_1^*Z_1=Z_2^*Z_2\), the map \(Z_1v\mapsto Z_2v\) is an isometry
between their images and extends to a unitary map of \(\mathbb C^3\).
Its action on the orthogonal complement can be adjusted to make its
determinant one.  In rank two the complement is one-dimensional and
the adjustment is unique, hence the stabilizer is trivial.  In rank
one the stabilizer is \(SU(2)\) on the two-dimensional complement.
At zero it is all of \(SU(3)\).  Surjectivity has the retained section
\[
 Z_p=\alpha^{-1/2}
 \begin{pmatrix}p^{1/2}\\0\quad0\end{pmatrix}.
\]
Compactness of \(SU(3)\) and continuity of the positive square root
give the quotient homeomorphism.  On \(\Ppos\) the section is smooth.
On \(\Nfuture\), the rank-one identity
\[
 p^{1/2}=\frac{p}{\sqrt{\operatorname{tr}p}}
\]
shows directly that the same section is smooth; together with the
constant orbit type \(SU(3)/SU(2)\), this gives the stated smooth
nonprincipal stratum.
\end{proof}

\begin{lemma}[Exact ranks of the Gram differential]
\label{wpr:lem:ranks}
\begin{enumerate}[label=\textup{(\roman*)}]
\item If \(\rankC Z=2\), then
\[
 \operatorname{im}D P_{\alpha,Z}=\J,\qquad
 \ker D P_{\alpha,Z}=T_Z(SU(3)\cdot Z),
\]
of real dimensions \(4\) and \(8\).
\item If \(\rankC Z=1\) and \(p=P_\alpha(Z)\), then
\[
 \operatorname{im}D P_{\alpha,Z}=p^{\perp_L}=T_p\Nfuture
\]
has real dimension \(3\).  The full ambient kernel has dimension \(9\),
but
\[
 \ker(D P_{\alpha,Z}|_{T_Z{\V}_1})=T_Z(SU(3)\cdot Z)
\]
has dimension \(5\), where \({\V}_1=\{Z:\rankC Z=1\}\).  The remaining
four real kernel directions are
\[
 \operatorname{Hom}_{\mathbb C}
 \bigl(\ker Z,(\operatorname{im}Z)^\perp\bigr);
\]
they raise rank and are not colour gauge.
\item At \(Z=0\), \(D P_{\alpha,0}=0\).
\end{enumerate}
\end{lemma}
\begin{proof}
For rank two and arbitrary \(A=A^*\), retain
\[
 \delta Z_A=\frac1{2\alpha}Z(Z^*Z)^{-1}A.
\]
Substitution into \eqref{wpr:eq:DP} gives
\(D P_{\alpha,Z}(\delta Z_A)=A\).  Hence the kernel has dimension
\(12-4=8\), equal to the free colour-orbit tangent.

For rank one write, without fixing either norm,
\[
 Z=ab^*,\qquad a\in\mathbb C^3\setminus\{0\},\quad
 b\in\mathbb C^2\setminus\{0\},\quad r=a^*a>0.
\]
Then \(p=\alpha rbb^*\).  If \(w=\delta Z^*a\), every
\(w\in\mathbb C^2\) occurs by taking \(\delta Z=aw^*/r\), and
\[
 \operatorname{im}D P_{\alpha,Z}
 =\{\alpha(wb^*+bw^*):w\in\mathbb C^2\}.
\]
The kernel of \(w\mapsto wb^*+bw^*\) is
\(\{itb:t\in\mathbb R\}\), so the image has dimension three.  It is
contained in \(p^{\perp_L}\), and equality follows by dimension.

Use the exact orthogonal splittings
\[
 \mathbb C^2=(\ker Z)^\perp\oplus\ker Z,\qquad
 \mathbb C^3=\operatorname{im}Z\oplus(\operatorname{im}Z)^\perp.
\]
Retaining the nonzero singular value \(\sigma>0\), adapted unit vectors
give
\[
 Z=\begin{pmatrix}\sigma&0\\0&0\\0&0\end{pmatrix},\qquad
 \delta Z=\begin{pmatrix}a_0&b_0\\c_1&d_1\\c_2&d_2\end{pmatrix},
\]
and
\[
 \alpha^{-1}D P_{\alpha,Z}(\delta Z)=
 \begin{pmatrix}
 2\sigma\operatorname{Re}a_0&\sigma b_0\\
 \sigma\overline{b_0}&0
 \end{pmatrix}.
\]
The ambient kernel has \(\operatorname{Re}a_0=b_0=0\) and dimension
nine.  The linearized rank-one equations also impose \(d_1=d_2=0\).
The five remaining directions are the \(SU(3)/SU(2)\)-orbit tangent;
the \(d_1,d_2\) block is the displayed complex Hom-space.  The zero
case follows from \eqref{wpr:eq:DP}.
\end{proof}

\subsection{Forms after colour reduction}
Define
\[
 \overline{\mathcal E}
 ={\J}_x\times{\Psd}_p\times\mathbb R_e
 \cong\mathcal E/SU(3),\qquad
 q(x,Z,e)=(x,P_\alpha(Z),e).
\]

\begin{proposition}[Reduced one-form and two-form]
\label{wpr:prop:reduced-forms}
The forms are basic for the colour action and satisfy
\[
 \theta=q^*\overline\theta,\qquad
 \Omega=q^*\overline\Omega,
\]
where
\begin{align}
 \overline\theta&=\langle p,\dd x\rangle_L,\label{wpr:eq:bartheta}\\
 \overline\Omega&=-\dd\overline\theta,\label{wpr:eq:barOmega}\\
 \overline\Omega((v_x,v_p),(w_x,w_p))
 &=\langle v_x,w_p\rangle_L-\langle w_x,v_p\rangle_L.
 \label{wpr:eq:barOmega-explicit}
\end{align}
After omitting the multiplier coordinate, this is symplectic on
\(\J\times\Ppos\).  On \(\J\times\Nfuture\), its kernel is exactly
\[
 \ker\overline\Omega_{(x,p)}=\SpanR\{(p,0)\}.
\]
On \(\J\times\{0\}\), the stratum-restricted two-form is zero.
\end{proposition}
\begin{proof}
Basicness follows from Lemma~\ref{wpr:lem:null}; the formulas follow by
differentiation.  The rank-two assertion uses that \(\Ppos\) is open in
\(\J\).  At nonzero rank-one \(p\), one has
\[
 T_p\Nfuture=p^{\perp_L}.
\]
If \((v_x,v_p)\) annihilates \eqref{wpr:eq:barOmega-explicit}, testing
arbitrary \(w_x\) gives \(v_p=0\), and testing
\(w_p\in p^{\perp_L}\) gives
\(v_x\in(p^{\perp_L})^{\perp_L}=\SpanR\{p\}\).
At \(p=0\), the zero stratum has no momentum tangent direction.
\end{proof}

\subsection{Positive mass}
For \(m_0>0\), define the retained future shell
\[
 \mathcal H^+_{m_0,c}
 =\{p\in\Ppos:\det p=m_0^2c^2\}.
\]
The parameter \(\alpha\) remains present in the exact quotient map
\(p=\alpha Z^*Z\); the displayed subset of momentum space itself
depends on \(m_0,c\).

\begin{theorem}[Exact massive reduction]
\label{wpr:thm:massive}
Let \(m_0>0\).  Then:
\begin{enumerate}[label=\textup{(\alph*)}]
\item The constraint surface
\[
 \Sigma_{m_0,\alpha,c}
 =\J\times\{Z:\det P_\alpha(Z)=m_0^2c^2\}\times\mathbb R_e
\]
is smooth and the colour action on it is free.
\item Its colour quotient is
\[
 \overline\Sigma_{m_0,c}
 =\J\times\mathcal H^+_{m_0,c}\times\mathbb R_e.
\]
Before colour reduction the kernel of the restricted two-form is the
direct sum of the eight-dimensional colour-orbit tangent,
\(\SpanR\{(p,0)\}\), and \(\SpanR\{\partial_e\}\).  After colour
reduction and omission of \(e\), only \(\SpanR\{(p,0)\}\) remains.
\item The equations are
\[
 \dot p=0,\qquad \dot x=e\,p,\qquad
 \det p=m_0^2c^2,
\]
while \(\dot e\) and motion in the colour fibre are arbitrary.
\item The physical phase space
\[
 \mathcal R_{m_0,c}
 =
 \frac{\J\times\mathcal H^+_{m_0,c}}
 {(x,p)\sim(x+s p,p),\ s\in\mathbb R}
\]
is canonically symplectomorphic to \(T^*\mathcal H^+_{m_0,c}\).
\item The colour-reduced positive-mass model is exactly the standard
future-directed spinless first-order relativistic particle of mass
\(m_0\) and speed constant \(c\).
\end{enumerate}
\end{theorem}
\begin{proof}
The constraint implies \(\rankC Z=2\).  By
Lemma~\ref{wpr:lem:ranks}, \(P_\alpha\) is a submersion there, and
\[
 D(\det)_p=2\langle p,\mathord{\cdot}\rangle_L\neq0.
\]
This proves smoothness.  Theorem~\ref{wpr:thm:quotient} gives freeness
and the colour quotient.  Since
\[
 T_p\mathcal H^+_{m_0,c}=p^{\perp_L},
\]
the proof of Proposition~\ref{wpr:prop:reduced-forms} gives the
characteristic line.  Pullback through the Gram submersion adds exactly
the colour-orbit kernel, and the multiplier adds \(\partial_e\).

Surjectivity of \(D P_{\alpha,Z}\) turns \eqref{wpr:eq:EL3} into
\(\dot x=e p\); the other equations are \eqref{wpr:eq:EL1} and
\eqref{wpr:eq:EL2}.  A path with constant \(p\) lies in one free Gram
fibre and differs locally from a constant lift by a
\(\tau\)-dependent colour transformation.

For the physical identification, if
\(A\in T_p\mathcal H^+_{m_0,c}\), define
\[
 \xi_{[x],p}(A)=-\langle x,A\rangle_L.
\]
This is unchanged by \(x\mapsto x+s p\).  Conversely, the annihilator
of \(p^{\perp_L}\) is \(\SpanR\{p\}\), so every cotangent covector has
one and only one class \([x]\).  Hence
\[
 ([x],p)\longmapsto(p,\xi_{[x],p})
\]
is a fibrewise linear bijection to \(T^*\mathcal H^+_{m_0,c}\).
The characteristic \(\mathbb R\)-action is free because every shell
momentum is nonzero.  It is proper: choose any Euclidean norm on
\(\J\).  For compact sets \(K_1,K_2\subset
\J\times\mathcal H^+_{m_0,c}\), the momenta occurring in them have a
positive minimum Euclidean norm \(\delta\); if
\((x,p)\in K_1\) and \((x+s p,p)\in K_2\), boundedness of the two
position projections gives \(|s|\delta\leq C(K_1,K_2)\).
Thus the orbit space is smooth.  The smooth vector-bundle map
\[
 (x,p)\longmapsto
 \bigl(p,\;A\mapsto-\langle x,A\rangle_L\bigr)
\]
is fibrewise onto with kernel \(\SpanR\{p\}\), and therefore descends
with a smooth inverse.  The displayed bijection is consequently a
diffeomorphism.

The basic one-form on the shell is
\begin{equation}
 \theta_{\mathrm{basic}}=-\langle x,\dd p\rangle_L.
 \label{wpr:eq:basic}
\end{equation}
It is the pullback of the tautological cotangent one-form, and
\begin{equation}
 \overline\theta
 =\dd\langle p,x\rangle_L+\theta_{\mathrm{basic}}.
 \label{wpr:eq:boundary}
\end{equation}
Therefore
\(\overline\Omega=-\dd\theta_{\mathrm{basic}}\), the canonical
cotangent symplectic form under the convention
\(\omega_{\mathrm{can}}=-\dd\lambda_{\mathrm{can}}\).
The original \(\overline\theta\) is not horizontal on the massive
characteristic because \(\overline\theta(p,0)=m_0^2c^2\); the exact
boundary term in \eqref{wpr:eq:boundary} produces the reduced
one-form.

Finally the reduced action is
\[
 \int_I\left(\langle p,\dot x\rangle_L
 -\frac e2(\det p-m_0^2c^2)\right)\dd\tau.
\]
Because \(\Ppos\) is open in \(\J\), \(p\) has every infinitesimal
Hermitian variation.  This is the standard covariant first-order
spinless particle action on the future-timelike sector, with no
residual internal symplectic factor.
\end{proof}

\begin{proposition}[Exact Hamiltonian gauge quasi-symmetry]
\label{wpr:prop:gauge}
For a smooth real \(\varepsilon(\tau)\), set
\[
 \delta_\varepsilon x=\varepsilon P_\alpha(Z),\qquad
 \delta_\varepsilon Z=0,\qquad
 \delta_\varepsilon e=\dot\varepsilon.
\]
Then
\begin{equation}
 \delta_\varepsilon L
 =
 \frac{\dd}{\dd\tau}
 \left[
 \frac{\varepsilon}{2}
 \bigl(\det P_\alpha(Z)+m_0^2c^2\bigr)
 \right].
 \label{wpr:eq:gauge-boundary}
\end{equation}
For compactly supported \(\varepsilon\) this is an exact gauge
quasi-symmetry, whose phase-space part generates the shell
characteristic where that characteristic is regular.
\end{proposition}
\begin{proof}
Put \(p=P_\alpha(Z)\) and \(q=\det p=\langle p,p\rangle_L\).  Since
\(\dot q=2\langle p,\dot p\rangle_L\),
\begin{align*}
 \delta_\varepsilon L
 &=
 \langle p,\dot\varepsilon p+\varepsilon\dot p\rangle_L
 -\frac{\dot\varepsilon}{2}(q-m_0^2c^2)\\
 &=
 \frac{\dot\varepsilon}{2}(q+m_0^2c^2)
 +\frac{\varepsilon}{2}\dot q,
\end{align*}
which is \eqref{wpr:eq:gauge-boundary}.
\end{proof}


\subsection{Massless rank one}
Set
\[
 \mathcal C_{0,\alpha,c}
 =\{Z:\det P_\alpha(Z)=0\}
 =\{Z:\rankC Z\leq1\}.
\]
The constant \(c>0\) is retained in the model, although it cancels
from the set-theoretic \(m_0=0\) equation.

\begin{lemma}[The pulled-back massless constraint is irregular]
\label{wpr:lem:irregular}
Writing the rows of \(Z\) as \(Z_{a\bullet}\),
\begin{equation}
 \det P_\alpha(Z)=
 \alpha^2\sum_{1\leq a<b\leq3}
 |Z_{a1}Z_{b2}-Z_{a2}Z_{b1}|^2.
 \label{wpr:eq:CB}
\end{equation}
Consequently,
\[
 \dd(\det P_\alpha)_Z=0
 \quad\text{for every }Z\in\mathcal C_{0,\alpha,c}.
\]
The single real multiplier equation describes the rank-at-most-one
set, but its differential supplies none of the four real normal
equations on the smooth rank-one stratum.
\end{lemma}
\begin{proof}
Equation \eqref{wpr:eq:CB} is Cauchy--Binet with the factor
\(\alpha^2\) retained.  On its zero set every complex minor vanishes,
so differentiation of the sum of squared absolute values gives zero.
Equivalently, at rank one,
\[
 \dd(\det P_\alpha)_Z(\delta Z)
 =2\langle p,D P_{\alpha,Z}(\delta Z)\rangle_L=0
\]
by Lemma~\ref{wpr:lem:ranks}; at \(Z=0\) use
\(D P_{\alpha,0}=0\).
\end{proof}

\begin{theorem}[Nonzero massless equations and physical quotient]
\label{wpr:thm:massless}
Suppose \(m_0=0\) and \(\rankC Z=1\), and put
\(p=P_\alpha(Z)\in\Nfuture\).  Then:
\begin{enumerate}[label=\textup{(\alph*)}]
\item The colour fibre is \(SU(3)/SU(2)\), with stabilizer \(SU(2)\).
\item The complete equations are
\begin{equation}
 \dot p=0,\qquad
 \dot x=\lambda(\tau)p
 \quad\text{for some real }\lambda(\tau),\qquad
 e(\tau)\text{ arbitrary}.
 \label{wpr:eq:null-EL}
\end{equation}
No equation identifies \(\lambda\) with the written multiplier \(e\).
\item On the intrinsic smooth rank-one shell, the kernel before
colour reduction is the direct sum of
\[
 T_Z(SU(3)\cdot Z),\qquad
 \SpanR\{(p,0)\},\qquad
 \SpanR\{\partial_e\}.
\]
After colour reduction and omission of \(e\), only the characteristic
\(\SpanR\{(p,0)\}\) remains.
\item The physical phase space
\[
 \mathcal R_0^\times
 =
 \frac{\J\times\Nfuture}
 {(x,p)\sim(x+s p,p)}
\]
is canonically symplectomorphic to \(T^*\Nfuture\), with reduced
one-form \(-\langle x,\dd p\rangle_L\) and its canonical two-form.
No quotient by positive rescaling of \(p\) occurs.
\item Thus the unparameterized spacetime curves and physical
symplectic quotient agree with the standard future-directed nonzero
spinless massless particle.  The original triples \((x,Z,e)\), their
off-shell neighbourhood, and their multiplier equations are not
exactly the standard massless system.
\end{enumerate}
\end{theorem}
\begin{proof}
The fibre statement is Theorem~\ref{wpr:thm:quotient}.  By
Lemma~\ref{wpr:lem:ranks}, equation \eqref{wpr:eq:EL3} says
\[
 \dot x-e p\in(p^{\perp_L})^{\perp_L}=\SpanR\{p\}.
\]
Since \(ep\) lies in this line, this is equivalent to
\(\dot x=\lambda p\), with no relation between \(\lambda\) and \(e\).
Equation \eqref{wpr:eq:EL2} gives \(\dot p=0\).  A path in a fixed
rank-one Gram fibre differs locally from a constant representative by
a \(\tau\)-dependent \(SU(3)\) transformation; the \(SU(2)\)
stabilizer makes this gauge description reducible.

On the intrinsic rank-one manifold,
Lemma~\ref{wpr:lem:ranks} identifies the vertical kernel with the
five colour directions.  Proposition~\ref{wpr:prop:reduced-forms}
adds the one spacetime characteristic, and \(e\) adds
\(\partial_e\).

Since \(T_p\Nfuture=p^{\perp_L}\), the massive cotangent construction
applies without division by \(\det p\):
\[
 ([x],p)\longmapsto
 \left(p,\;A\longmapsto-\langle x,A\rangle_L\right).
\]
Every \(p\in\Nfuture\) is nonzero.  The freeness, properness, and smooth
vector-bundle argument in the massive proof applies to compact subsets
of \(\J\times\Nfuture\) verbatim, and proves that this map is a
diffeomorphism to \(T^*\Nfuture\).  Moreover,
\eqref{wpr:eq:boundary} identifies the canonical form.  Reparametrizing
the worldline changes \(x\) along \(p\); it does not identify distinct
positive magnitudes of \(p\).  Lemma~\ref{wpr:lem:irregular} and the
absence of \(\lambda=e\) establish the final inequivalence statement.
\end{proof}

\begin{counterexample}[Rank-one multiplier mismatch]
\label{wpr:ce:rank1}
Retain arbitrary \(s>0\), and set
\[
 Z_s=\begin{pmatrix}s&0\\0&0\\0&0\end{pmatrix},\qquad
 p_s=\alpha\begin{pmatrix}s^2&0\\0&0\end{pmatrix}.
\]
The path
\[
 Z(\tau)=Z_s,\qquad x(\tau)=\tau p_s,\qquad e(\tau)=0
\]
satisfies all equations in Proposition~\ref{wpr:prop:EL} for
\(m_0=0\).  The standard action with momentum independently varied in
\(\J\) and the same \(e=0\) would require \(\dot x=e p_s=0\).

Under Proposition~\ref{wpr:prop:gauge}, a solution
\(\dot x=\lambda p\), \(\dot p=0\), transforms by
\[
 \lambda\longmapsto\lambda+\dot\varepsilon,\qquad
 e\longmapsto e+\dot\varepsilon.
\]
Thus \(\lambda-e\) is invariant under the stated Hamiltonian gauge;
the discrepancy remains when the original multiplier is retained.
\end{counterexample}

\begin{counterexample}[Four non-colour first-order null directions]
\label{wpr:ce:transverse}
Retain arbitrary \(s>0\) and \(t\in\mathbb C\setminus\{0\}\), and put
\[
 Z_s=\begin{pmatrix}s&0\\0&0\\0&0\end{pmatrix},\qquad
 \Delta_t=\begin{pmatrix}0&0\\0&t\\0&0\end{pmatrix}.
\]
Then
\begin{equation}
 D P_{\alpha,Z_s}(\Delta_t)=0,\qquad
 \det P_\alpha(Z_s+\varepsilon\Delta_t)
 =\alpha^2s^2|t|^2\varepsilon^2
 \quad(\varepsilon\in\mathbb R).
 \label{wpr:eq:transverse}
\end{equation}
The vector \(\Delta_t\) is not tangent to a colour orbit, since every
\(\xi Z_s\) has zero second column.  It is not tangent to the rank-one
variety, because the perturbed matrix has rank two for
\(\varepsilon\neq0\).  Allowing both complex coordinates in
\((\operatorname{im}Z_s)^\perp\) gives the four real transverse
directions of Lemma~\ref{wpr:lem:ranks}.  They are second-order
shell-opening directions, not gauge.
\end{counterexample}

\subsection{The zero-momentum vertex}
\begin{theorem}[Additional singular branch]
\label{wpr:thm:zero}
Let \(m_0=0\) and \(Z=0\).  Then \(p=0\), the stabilizer is all of
\(SU(3)\), and both \(\theta\) and \(\Omega\) vanish on
\(\J\times\{0\}\times\mathbb R_e\).  Every pair of paths
\[
 x:I\to\J,\qquad e:I\to\mathbb R,
\]
with \(Z(\tau)\equiv0\), satisfies the Euler--Lagrange equations.

The phase part of the Hamiltonian gauge vanishes at \(p=0\), while the
presymplectic kernel contains every \(x\)-direction.  Quotienting by
the actual colour and Hamiltonian-gauge orbits therefore leaves the
whole \(\J\) branch, whereas quotienting by all characteristic leaves
would collapse it to one point.  This is an exact failure of regular
Dirac reduction at the vertex.
\end{theorem}
\begin{proof}
At \(Z=0\), one has \(p=0\) and \(D P_{\alpha,0}=0\).  Every equation
in Proposition~\ref{wpr:prop:EL} is identically satisfied, independent
of \(\dot x,e,\dot e\).  The form and gauge statements follow from
\eqref{wpr:eq:forms} and
\(\delta_\varepsilon x=\varepsilon p=0\).
\end{proof}

\begin{counterexample}[Stationary zero momentum with nonconstant
position]
\label{wpr:ce:zero}
The retained paths
\[
 x(\tau)=\begin{pmatrix}\tau^2&0\\0&0\end{pmatrix},\qquad
 Z(\tau)=0,\qquad e(\tau)=\sin\tau
\]
are stationary for the present \(m_0=0\) action.  The standard
independent-momentum action gives \(\dot x=e p=0\) at \(p=0\), and
therefore excludes this path.  Removing \(p=0\), as is conventional
for physical massless phase space, removes this vertex counterexample
but does not regularize the off-shell rank-one constraint.
\end{counterexample}

\subsection{Kernel dimensions and exact verdict}
\begin{proposition}[Every kernel jump]
\label{wpr:prop:kernels}
Omit the always-null multiplier direction for the table.  On the
ambient \(16\)-real-dimensional \({\J}_x\times{\V}_Z\),
\begin{center}
\begin{tabular}{c c c c}
\toprule
\(\rankC Z\)&\(\operatorname{rank}_{\mathbb R}D P_{\alpha,Z}\)&
\(\operatorname{rank}_{\mathbb R}\Omega\)&
\(\dim_{\mathbb R}\ker\Omega\)\\
\midrule
\(2\)&\(4\)&\(8\)&\(8\)\\
\(1\)&\(3\)&\(6\)&\(10\)\\
\(0\)&\(0\)&\(0\)&\(16\)\\
\bottomrule
\end{tabular}
\end{center}
On the \(12\)-real-dimensional intrinsic rank-one stratum
\({\J}_x\times{\V}_1\), the pulled-back form has rank \(6\) and nullity
\(6\): five colour directions and one spacetime characteristic.
\end{proposition}
\begin{proof}
Using the Lorentz pairing to identify \(\J\cong\J^*\), the matrix of
\(\Omega\) is
\[
 \begin{pmatrix}
 0&G_L D P_{\alpha,Z}\\
 -(D P_{\alpha,Z})^{\mathsf T}G_L&0
 \end{pmatrix}.
\]
Its rank is twice the rank of \(D P_{\alpha,Z}\).  Apply
Lemma~\ref{wpr:lem:ranks} in the ambient space and on \(T_Z{\V}_1\).
\end{proof}

\begin{theorem}[Final comparison]
\label{wpr:thm:verdict}
The strongest correct result is:
\begin{enumerate}[label=\textup{(\arabic*)}]
\item for \(m_0>0\), colour reduction gives exactly the standard
future-timelike first-order spinless particle, including its equations,
constraint characteristic, and physical cotangent symplectic space;
\item for \(m_0=0\) and \(p\neq0\), colour and characteristic reduction
gives exactly the standard future-null physical phase space and
unparameterized null trajectories;
\item the full \(m_0=0\) pulled-back multiplier action is not the
standard massless variational system: its constraint has zero
differential, its rank-one equation is \(\dot x=\lambda p\) rather
than \(\dot x=e p\), its ambient kernel contains four non-colour
transverse directions, and its vertex admits arbitrary \(x(\tau)\).
\end{enumerate}
The exact relating morphism is \(P_\alpha(Z)=\alpha Z^*Z\); the exact
failure locus is the singular differential of
\(\det\circ P_\alpha\) on the rank-at-most-one variety.
\end{theorem}
\begin{proof}
Combine Theorem~\ref{wpr:thm:massive},
Theorem~\ref{wpr:thm:massless},
Lemma~\ref{wpr:lem:irregular}, and
Counterexamples~\ref{wpr:ce:rank1}--\ref{wpr:ce:zero}.
\end{proof}

\subsection{Deterministic verification companion}
The companion program
\texttt{verify\_worldline\_reduction.py} checks the determinant
polarization, the rank-two Gram right inverse, all Gram and
presymplectic ranks, all colour-orbit ranks, the intrinsic rank-one
tangent calculation, the zero gradient and quadratic transverse
opening of the massless constraint, the cotangent-form sign, and
\eqref{wpr:eq:gauge-boundary}, using exact symbolic arithmetic.

\subsection{Comparison with the extended exact Lorentz packet}
\label{wpr:app:packet}

This subsection records a bounded line-specific audit of source lines
875--1654.  The audited source file is
\begin{center}
\small\path{exact_lorentz_observation_bridge.tex}.
\end{center}
It is stored beneath
\path{supporting_materials/workbench/evidence} in the directory
\begin{center}
\small\path{worldline_presymplectic_reduction_20260904}.
\end{center}
Its file SHA-256 is
\[
\texttt{2BE16E1DC71B2E5DC04238F0FBE440138258D0825A4841AB2F29662587681599}.
\]
That packet adds a third column \(\phi\in\mathbb C^3\), writes
\(B=[Z\ \phi]\), and uses
\[
 G=Z^*Z,\quad z=Z^*\phi,\quad
 \lambda=\phi^*\phi,\quad \omega=\det B.
\]
Its mass convention is \(m^2_{\mathrm{packet}}=\det p\), with physical
mass \(M_{\mathrm{phys}}=\sqrt{\det p}/c\).  The convention in the
audited \(S^6\) module is
\(m_{\alpha,c}^2=\det p/c^2\).  The exact dictionary is therefore
\[
 m^2_{\mathrm{packet}}=c^2m_{\alpha,c}^2,\qquad
 M_{\mathrm{phys}}=m_{\alpha,c};
\]
there is no mathematical disagreement.

\subsubsection{Claim ledger}

\begin{center}
\begin{tabular}{>{\raggedright\arraybackslash}p{0.15\textwidth}
                >{\raggedright\arraybackslash}p{0.20\textwidth}
                >{\raggedright\arraybackslash}p{0.55\textwidth}}
\toprule
Source lines&Verdict&Exact scope\\
\midrule
875--950&validated&
\(p=\alpha Z^*Z\) is a proper surjection, its fibres are the
\(SU(3)\)-orbits, and
\(\det p=\alpha^2\sum_{a<b}|\Delta_{ab}(Z)|^2\) has range
\([0,\infty)\).  The zero fibre is rank at most one.\\
\addlinespace
952--1027&validated as a multiplication-operator model&
With the stated pushforward measure, the observer-energy and
mass-squared multiplication operators have essential ranges
\([0,\infty)\), no point spectrum, and no positive gap after a vacuum
is adjoined.  This is a kinematic/free spectral statement, not a
statement about an interacting field Hamiltonian.\\
\addlinespace
1029--1476&validated&
For \(\lambda_0>0\), \(\omega_0\neq0\), the constrained quotient,
associated-ball description, momentum fibres, and sharp lower bound
\(\det p\geq\alpha^2|\omega_0|^2/\lambda_0\) are algebraic and
differential-topological statements.  Since \(\omega_0\neq0\) forces
\(B\) invertible and \(Z\) rank two, the massless rank-one defect
proved above does not invalidate these claims.\\
\addlinespace
1482--1531&valid invariance and constraints, incomplete dynamics&
The displayed action is invariant and its multiplier and \(x\)
variations give the stated constraints and \(\dot p=0\).  Its omitted
\((Z,\phi)\) equations materially change the dynamical interpretation;
Theorem~\ref{wpr:thm:packet-action} below computes them intrinsically.\\
\addlinespace
1533--1603&validated with scope qualification&
The multiplication spectra and the direct-sum/Fock vacuum gap equal
the displayed threshold.  The threshold is inserted by fixing
\(\lambda=\lambda_0\) and \(\omega=\omega_0\neq0\); it is not a
dynamically generated interacting mass gap and is not a Yang--Mills
mass-gap theorem.\\
\addlinespace
1605--1633&validated and scope-sharpening&
The two explicit degenerating sequences prove that controlling only
one of \(\lambda\) and \(\omega\) gives no positive threshold.  Thus
the positive lower endpoint is exactly constraint-selected.\\
\addlinespace
1635--1653&validated&
The diagonal boost excludes a positive-definite quadratic form on the
spinor carrier invariant under all of \(SL_2(\mathbb C)\).\\
\bottomrule
\end{tabular}
\end{center}

The independent replay and line-by-line claim ledger for this packet are
retained under the project-relative path
\path{supporting_materials/workbench/evidence/}\allowbreak
\path{exact_lorentz_observation_bridge_audit_20260904}.
The packet checker passes \(33/33\) assertions and the separately written
checker passes \(30/30\).  The audit-report SHA--256 is
\texttt{c8d01af28836340798c2fa7e27f8e097\allowbreak%
d67e0b8d7c13f75271af473958d391df}; the promotion manifest, which verifies
all retained audit artifacts, has SHA--256
\par\noindent
\texttt{52f53746558fa219260b224abdb71c86d2cd6c13192774d0d2fb1a1ed71bcf35}.

\subsubsection{The full dynamics of the packet's constrained action}

Retain the packet's constants
\[
 \lambda_0>0,\qquad\omega_0\in\mathbb C^*,\qquad\alpha>0
\]
and its smooth colour quotient
\[
 \mathcal Q_{\lambda_0,\omega_0}
 \cong
 \operatorname{Herm}_2^0(\mathbb C)
 \times\{w\in\mathbb C^2:\|w\|^2<\lambda_0\}.
\]
Put
\begin{equation}
 \mu(w)=\frac{|\omega_0|}
 {\sqrt{\lambda_0-\|w\|^2}},\qquad
 G(A,w)=\mu(w)e^A,\qquad
 p(A,w)=\alpha G(A,w).
 \label{wpr:eq:packet-p}
\end{equation}
Eliminating the three regular multiplier constraints from source lines
1484--1495 gives the intrinsic action
\begin{equation}
 S_{\mathrm{packet}}[x,A,w]
 =\int_I\langle p(A,w),\dot x\rangle_L\,\dd\tau
 \label{wpr:eq:packet-action}
\end{equation}
on \({\J}_x\times\mathcal Q_{\lambda_0,\omega_0}\).  The use of the
intrinsic constrained action loses no Euler--Lagrange equation:
the constraint map is a submersion on \(\omega_0\neq0\), so the
Lagrange-multiplier theorem identifies ambient stationarity with
stationarity against every tangent vector of the constrained
manifold.  This is a statement about classical stationary paths; it
does not assert equality of quantum measures or path integrals after
formally integrating out the multipliers.

\begin{theorem}[Complete reduced dynamics of the constrained packet]
\label{wpr:thm:packet-action}
For the momentum projection in \eqref{wpr:eq:packet-p}:
\begin{enumerate}[label=\textup{(\roman*)}]
\item If \(w\neq0\), then \(D p_{(A,w)}\) has real rank \(4\), and its
kernel has dimension \(3\), tangent to the \(S^3\) momentum fibre.
The Euler--Lagrange equations of \eqref{wpr:eq:packet-action} are
\[
 \dot p=0,\qquad \dot x=0,
\]
with motion along the fixed \(S^3\) fibre presymplectically null.
\item If \(w=0\), then
\[
 \det p=\frac{\alpha^2|\omega_0|^2}{\lambda_0},
\]
\(D p_{(A,0)}\) has rank \(3\), its image is
\(p^{\perp_L}\), and its kernel consists of all four real
\(w\)-directions.  The pointwise equations are
\[
 \dot p=0,\qquad \dot x=\lambda(\tau)p.
\]
Every actual solution remaining on the threshold has constant
\((A,w)\), because the momentum fibre over a threshold momentum is a
single point.
\item Consequently, the source action is not a family of ordinary
massive particle actions over all masses above the lower endpoint.
Only its threshold stratum has the ordinary worldline characteristic;
strict-interior configurations have \(\dot x=0\).  The observer
multiplication operator of source lines 1533--1603 is a separate
positive-energy representation and is not the Hamiltonian generating
\(\tau\)-evolution in \eqref{wpr:eq:packet-action}.
\end{enumerate}
\end{theorem}

\begin{proof}
For fixed \(w\), varying \(A\) varies \(p\) through the
three-dimensional fixed-determinant future hyperboloid, so the image
of the \(A\)-differential is \(p^{\perp_L}\).  Direct differentiation
without suppressing any constant gives
\begin{equation}
 D\mu_w(\delta w)
 =
 \frac{|\omega_0|\operatorname{Re}\langle w,\delta w\rangle}
 {(\lambda_0-\|w\|^2)^{3/2}}.
 \label{wpr:eq:Dmu}
\end{equation}
When \(w\neq0\), a radial \(\delta w\) makes
\eqref{wpr:eq:Dmu} nonzero and changes \(\det p=\alpha^2\mu(w)^2\).
This supplies a fourth direction transverse to \(p^{\perp_L}\).
The remaining \(w\)-directions satisfy
\(\operatorname{Re}\langle w,\delta w\rangle=0\), form a
three-dimensional space, and are exactly tangent to the \(S^3\)
momentum fibre.

At \(w=0\), equation \eqref{wpr:eq:Dmu} vanishes in every direction,
so the four \(w\)-directions lie in \(\ker Dp\) and the image is only
\(p^{\perp_L}\).  These directions open the mass above its minimum at
second order:
\[
 \det p(A,\varepsilon v)
 =
 \frac{\alpha^2|\omega_0|^2}
 {\lambda_0-\varepsilon^2\|v\|^2}.
\]

For an action \(\int\langle p(q),\dot x\rangle_L\dd\tau\), variation
of \(x\) gives \(\dot p=0\), while variation of \(q\) gives
\[
 \langle Dp_q(\delta q),\dot x\rangle_L=0
 \quad\text{for all }\delta q.
\]
If \(w\neq0\), the image is all of \(\J\), hence \(\dot x=0\).
If \(w=0\), the image is \(p^{\perp_L}\), hence
\(\dot x\in\SpanR\{p\}\).  The packet's own momentum-fibre theorem
states that the fibre over a threshold momentum is a singleton, which
completes the solution statement.
\end{proof}

\subsubsection{Relation to the rank-one irregularity}

The packet's fixed-\(\omega_0\neq0\) constrained manifold never meets
\(\rankC Z\leq1\), so the zero differential of
\(\det\circ P_\alpha\) does not invalidate its topology or positive
lower bound.  The exact relation is instead:
\begin{enumerate}
\item on the unrestricted carrier, the rank-one and zero strata give
the massless irregularities proved in
Theorems~\ref{wpr:thm:massless} and \ref{wpr:thm:zero};
\item on the packet's constrained carrier, \(Z\) remains rank two, but
the momentum projection develops a different four-real-dimensional
critical kernel at its positive-mass threshold \(w=0\);
\item the algebraic inequality is a sharp, externally
constraint-selected threshold; a dynamically generated gap would
require a specified interacting Hamiltonian or field operator whose
spectrum produces the scale rather than receiving
\(\lambda_0,\omega_0\) as fixed input.
\end{enumerate}

\endgroup
